% Tabela-síntese da seção 3.3 (média-variância aplicada ao mercado de criptoativos)
\SinteseInicio{Estudos de média-variância aplicados ao mercado de criptoativos}{tab:sintese_mv_cripto}

\SinteseAutor{brauneis2019cryptocurrency} &
Diversificação e desempenho da \acs{mv}. &
Universo de 500 criptomoedas entre 2015 e 2017; comparação com ativos isolados e carteira 1/N. &
Quase 98\% das correlações pareadas ficam entre \textminus0,10 e 0,20; fora da amostra, a carteira 1/N supera mais de 75\% das carteiras \acs{mv} avaliadas em Sharpe e equivalente de certeza. &
Oferece justificativa empírica para comparar \acs{adaboost} + \acs{mv} com 1/N, referência prevista na metodologia. \\
\SinteseLinha

\SinteseAutor{jeleskovic2024cryptocurrency} &
Dependência e otimização em criptoativos. &
Modelo \acs{garch}-cópula no arcabouço de Markowitz; carteiras com criptoativos e ativos tradicionais. &
Aponta benefícios em carteiras mistas e dificuldades para estimar dependências em um mercado de correlações instáveis. &
Contextualiza as limitações de estimar o risco de carteiras de criptoativos. \\
\SinteseLinha

\SinteseAutor{zouaoui2025portfolio} &
Concentração e limites de Markowitz. &
Comparação entre Markowitz e abordagem com \acs{lstm} para estimar pesos de carteiras de criptomoedas. &
Relata concentração em poucos ativos e Sharpe menos favorável na carteira Markowitz avaliada; a \acs{lstm} obtém melhor Sharpe no experimento. &
Acrescenta a concentração da carteira às limitações da \acs{mv}; o resultado com \acs{lstm} não valida diretamente \acs{adaboost} + \acs{mv}. \\

\SinteseFim
